% GENERATED FILE. Edit canonical lessons, then run npm run generate:books.
% canonical-source-hash: fnv1a64:98bf0d2bda9177c5
% canonical-lessons: TA-C25-iniya-iravu, TA-W10-read-naan, TA-C25-goodnight-alternatives

\chapter{Good Night}
\label{ch:ta-good-night}
\hlchaptermodality{25}

\begin{chapteropening}
\textbf{By the end of this chapter:} \emph{I can say good night in Tamil out of roots I already own, and trace every one of those forms back to \ta{நல்}.}

Say nalla iravu and iravu vaṇakkam, then run nal- forward into nalam, naṉṟi and nalla.
\end{chapteropening}

\section[\ta{இனிய} \ta{இரவு}]{\ta{இனிய} \ta{இரவு} (iṉiya iravu) — \textquotedblleft{}good night,\textquotedblright{} and Tamil breaking the pattern one more time}

\label{lesson:TA-C25-iniya-iravu}



\begin{quote}
Hindi, Kannada, Telugu, and Malayalam all reached for the same Sanskrit phrase to say \textquotedblleft{}good night.\textquotedblright{} This lesson closes the arc with the one that didn't — and it's a pattern you've actually seen from Tamil before, all the way back in Chapter 1.
\end{quote}

\begin{cousinweb}
\textbf{\ta{இனிய} \ta{இரவு}} (\textbf{iṉiya iravu}) — \textquotedblleft{}\textbf{good night}\textquotedblright{} — literally \textquotedblleft{}\textbf{sweet, pleasant night}.\textquotedblright{} \textbf{\ta{இனிய}} (\emph{iṉiya}, \textquotedblleft{}sweet, pleasant, dear\textquotedblright{}) is native \textbf{Dravidian}, from \textbf{Proto-Dravidian} \emph{\textbf{*in-}} (\textquotedblleft{}sweet\textquotedblright{}) — the same root behind \textbf{\ta{இனிப்பு}} (\emph{iṉippu}, \textquotedblleft{}sweet, dessert\textquotedblright{}). Paired with \textbf{\ta{இரவு}} (\emph{iravu}, \textquotedblleft{}night,\textquotedblright{} already met last lesson), the whole phrase is built from Tamil's own word-stock, start to finish.
\end{cousinweb}

\begin{culture}
Tamil's Dravidian sisters and Hindi all reached for the \textbf{same} Sanskrit phrase: Hindi's \textbf{\dv{शुभ} \dv{रात्रि}}, Kannada's \textbf{\ka{ಶುಭ} \ka{ರಾತ್ರಿ}}, Telugu's \textbf{\te{శుభ} \te{రాత్రి}}, Malayalam's \textbf{\ml{ശുഭ} \ml{രാത്രി}} — all four built from the identical tatsama pair, \textbf{śubh} (\textquotedblleft{}to be beautiful\textquotedblright{}) + \textbf{rātri} (\textquotedblleft{}night\textquotedblright{}). Tamil's standard phrase reaches for \textbf{none} of that — no śubha-descended \textquotedblleft{}good night\textquotedblright{} turned up as a commonly used Tamil form in the phrasebook sources checked for this lesson. Instead, Tamil built its own phrase from its own vocabulary.
\end{culture}

\subsection*{Your turn}
Say these aloud:

\begin{itemize}
\raggedright
  \item \textquotedblleft{}iṉiya iravu\textquotedblright{} — \textquotedblleft{}good night,\textquotedblright{} literally \textquotedblleft{}sweet night,\textquotedblright{} fully native
  \item the pattern it breaks — every other language this arc shares one Sanskrit śubh(a) rātri phrase; Tamil doesn't
\end{itemize}

\begin{tcolorbox}[breakable,colback=teal!4,colframe=teal!35!black,title={Before you move on}]
What does \textbf{\ta{இனிய} \ta{இரவு}} literally mean, and what root does \textbf{\ta{இனிய}} come from? (\textquotedblleft{}\textbf{Sweet/pleasant night}\textquotedblright{} — Proto-Dravidian \emph{\textbf{in-}}, \textquotedblleft{}sweet.\textquotedblright{}) Does Tamil's \textquotedblleft{}good night\textquotedblright{} share the same Sanskrit śubha rātri phrase as Hindi, Kannada, Telugu, and Malayalam? (\textbf{No} — no such form turned up in the sources checked for this lesson; Tamil's standard phrase is fully native, breaking the pattern shared by all four other languages.) Next: the \ta{ா} sign, and reading \ta{நான்} off the page.
\end{tcolorbox}
\section[nāṉ]{\ta{நான்} — the sign that stands after}

\label{lesson:TA-W10-read-naan}



\begin{quote}
Two places a vowel sign has sat so far: \textbf{\ta{ி}} goes above and to the right, \textbf{\ta{ை}} and \textbf{\ta{ெ}} go to the left. Here is the third place, and it is the one you would have guessed first.
\end{quote}

\subsection*{Script you'll notice: the \ta{ா} sign}
\textbf{\ta{ா}} is the sign for long \emph{ā}, and this book's script data describes it plainly: \textbf{a vertical stroke with a small top hook, written after the consonant.}

\begin{quote}
\textbf{\ta{ம} + \ta{ா} $\to$ \ta{மா}}, said \emph{mā}.
\end{quote}

That pairing is the example the data itself gives, on a letter you can already draw.

So the three positions, all seen now:

\noindent
\begin{tabularx}{\linewidth}{@{}>{\raggedright\arraybackslash}X>{\raggedright\arraybackslash}X>{\raggedright\arraybackslash}X@{}}
\toprule
\textbf{sign} & \textbf{sound} & \textbf{where it sits} \\
\midrule
\textbf{\ta{ா}} & \emph{ā} & \textbf{after} the consonant \\
\textbf{\ta{ி}} & \emph{i} & above and to the right \\
\textbf{\ta{ை}} & \emph{ai} & to the left, spoken after \\
\textbf{\ta{ெ}} & \emph{e} & to the left, spoken after \\
\bottomrule
\end{tabularx}

Notice also that \textbf{\ta{ஆ}}, the independent long \emph{ā} that opens a word, and \textbf{\ta{ா}}, the sign, carry the same lengthening. One is a letter standing on its own; the other hangs off a consonant.

\subsection*{Script you'll notice: build the word}
\noindent
\begin{tabularx}{\linewidth}{@{}>{\raggedright\arraybackslash}X>{\raggedright\arraybackslash}X>{\raggedright\arraybackslash}X@{}}
\toprule
\textbf{piece} & \textbf{} & \textbf{} \\
\midrule
\textbf{\ta{நா}} & \emph{nā} — \ta{ந} with the \ta{ா} sign after it & \ta{ந} drawn in \textbf{\ta{நன்றி}} \\
\textbf{\ta{ன்}} & bare \emph{ṉ} — \textbf{\ta{ன}} under a puḷḷi & \ta{ன} drawn there too, the alveolar \emph{n} \\
\bottomrule
\end{tabularx}

\begin{quote}
\textbf{\ta{நா} · \ta{ன்} $\to$ \ta{நான்}}
\end{quote}

Two pieces and one new mark. \textbf{\ta{நான்}} is \textquotedblleft{}I,\textquotedblright{} and its possessive is the \textbf{\ta{என்}} (\textquotedblleft{}my\textquotedblright{}) you read in chapter 21 — so \emph{nāṉ} and \emph{eṉ} now both come off the page.

\subsection*{Your turn}
\begin{itemize}
\raggedright
  \item \emph{Read it:} \textbf{\ta{மா}} — then \textbf{\ta{நா}} — then \textbf{\ta{நான்}}
  \item \emph{Say it:} \textquotedblleft{}nāṉ\textquotedblright{} — \textquotedblleft{}I\textquotedblright{}
  \item \emph{Contrast:} \textbf{\ta{நா}} · \textbf{\ta{றி}} · \textbf{\ta{பெ}} — a sign after, a sign above-right, a sign before
  \item \emph{Contrast:} \textbf{\ta{ஆம்}} and \textbf{\ta{நான்}} — the same long \emph{ā}, once as a letter and once as a sign
\end{itemize}

\begin{tcolorbox}[breakable,colback=teal!4,colframe=teal!35!black,title={Before you move on}]
Read \textbf{\ta{நான்}}. Where does \textbf{\ta{ா}} sit relative to its consonant? (\textbf{After it} — a vertical stroke with a small top hook.) Which of Tamil's three \emph{n}-letters closes the word, and what is the dot doing? (\textbf{\ta{ன}}, the alveolar one; the puḷḷi \textbf{removes its built-in \emph{a}}.) Next: good night, out of roots you already own.
\end{tcolorbox}
\section[\ta{நல்ல} \ta{இரவு} / \ta{இரவு} \ta{வணக்கம்}]{\ta{நல்ல} \ta{இரவு} and \ta{இரவு} \ta{வணக்கம்} — old roots recombined}

\label{lesson:TA-C25-goodnight-alternatives}



\begin{quote}
Tamil's native \textquotedblleft{}sweet night\textquotedblright{} repeats the lexical choice that opened this book: build greetings from its own roots rather than the shared Sanskrit formula.
\end{quote}

\subsection*{You'll want to know: The Chapter 1 echo}
\textbf{\ta{வணக்கம்}} comes from native \emph{vaṇaṅku}, \textquotedblleft{}bow,\textquotedblright{} while neighboring greetings use Sanskrit \emph{namas-}. At the other end of the day, \textbf{\ta{இனிய} \ta{இரவு}} is native while the same neighbors share Sanskrit \emph{śubha rātri}.

Tamil also keeps two transparent alternatives:

\begin{itemize}
\raggedright
  \item \textbf{\ta{நல்ல} \ta{இரவு}} — \textquotedblleft{}good night,\textquotedblright{} using \emph{nalla}, from the \textbf{nal-} \textquotedblleft{}good\textquotedblright{} root already inside \emph{nalam} and \emph{naṉṟi}
  \item \textbf{\ta{இரவு} \ta{வணக்கம்}} — \textquotedblleft{}night greetings,\textquotedblright{} especially for close friends/family, recombining two fully known words
\end{itemize}

The learner does not memorize isolated formulas; known roots generate options.

\subsection*{Your turn}
\begin{itemize}
\raggedright
  \item \emph{Say it:} \textquotedblleft{}nalla iravu\textquotedblright{} — good night
  \item \emph{Say it:} \textquotedblleft{}iravu vaṇakkam\textquotedblright{} — night greetings
  \item \emph{Trace it:} \emph{nal-} $\to$ \emph{nalam / naṉṟi / nalla}
\end{itemize}

\begin{tcolorbox}[breakable,colback=teal!4,colframe=teal!35!black,title={Before you move on}]
Which root builds \textbf{\ta{நல்ல}}? (\textbf{Nal-}, good.) What two known words make \textbf{\ta{இரவு} \ta{வணக்கம்}}? (\textbf{Night + greetings}.) What lexical pattern connects Chapter 1 and this chapter? (\textbf{Tamil keeps native greeting roots where neighbors share Sanskrit formulas}.)
\end{tcolorbox}
